\chapter{Indexing Powder Diffraction Patterns}\index{powder diffraction indexing}\label{Chap:indexing}

One of the more challenging aspects of powder diffraction is finding a unit cell that indexes the pattern. By indexing, I mean that the unit cell generates reflection positions that match most, or better yet all, of the diffraction peaks.  Autoindexing\index{autoindexing|textbf} is where computer software is used to search for a unit cell that will index the pattern. In contrast, indexing of single crystal diffraction patterns is a routine task that has been automated since the 1960's. Even with multiple twins present, modern software can facilely index single crystal data. Not so, alas for autoindexing powder diffraction. This, in theory, should not be a complex task on  modern computers. Noting that \textit{h}, \textit{k} and \textit{l} values have can only have limited values if one places reasonable limits on the lengths of unit cell constants, then there are a large number of possible indices that can be assigned, so it should not be that complex to search through all possibilities. However, in reality there may be mismatches between the peak positions we measure and the actual reflection positions: two reflections may be close enough that observed position of a peak matches a location between the two reflections. There may be impurities present, so some peaks may belong to a second phase. Sample displacement, particularly in a Bragg-Brentano diffractometer, can offset all peaks. The advent of high-resolution synchrotron diffraction has certainly aided the process of indexing powder data by providing data with very minimal sample displacement and very sharp peaks.  

Additional phases are not the only reason why your cell may not index all lines of the powder pattern. It is also possible that the cell you have found is a subcell of the actual lattice, so looking at conventional crystallographic subgroups (discussed in Chapter \ref{Chap:Subgroups}) might provide a related lattice that does index the pattern. Another type of supercell would be a superspace expansion to the lattice in a 4th dimension. This additional dimension is where there is lattice repetition in a direction that is an irrational\footnote{Note that if the lattice vector were a rational combination of the \textit{a}, \textit{b} and \textit{c} lengths then this would create a large but conventional superlattice.} combination of the three axes of your cell.  This is known as an incommensurate or modulated lattice. This extra direction is labeled as a 4th dimension, but note that this direction is actually occurring in three dimensional space, so this is labeled as 3+1 dimensional superspace. Such structures are not common, but are not unheard of either. Theory allows for two or even three irrational lattice vectors, so 3+2 and 3+3 dimensional superspace structures are possible, and there exist examples of these from single-crystal data, but it is unlikely that they would ever be fit with powder diffaction data.  GSAS-II does handle superspace groups of the 3+1 type, but not the 3+2 or 3+3 types. Note that the 230 conventional space groups, expand to 775, 3338 and 12,584 space groups, respectively for 3+1, 3+2 and 3+3 superspace. Should you need higher symmetry than 3+1, the JANA software suite should help.\index{software!JANA} Indexing of superspace groups, as well as indexing of magnetic lattices, requires finding an expansion vector, often referred to as a k-vector\index{k-vector}, which will not be covered in this book. Work is being done to add k-vector search tools to GSAS-II, so perhaps this will appear in a future edition. 

GSAS-II provides several tools for indexing powder diffraction patterns. These are accessed using a histogram's ``Unit Cells List'' data tree entry.  It should be noted that there are a number of computational methods that have been implemented for autoindexing. Sometimes one method will succeed with a list of peak positions where others fail. Should the algorithm in GSAS-II fail (which is an implementation of a published algorithm from Topas), you are recommended to also try versions of the classic autoindexing programs, TREOR, ITO and DICVOL .\index{software!TREOR}\index{software!ITO}\index{software!DICVOL} My rather aged CMPR\index{software!CMPR} program provides access to these programs, but there are probably now much better options, depending on your operating system. 

\section{Autoindexing}
The autoindexing problem requires the most precise possible peak positions, which are obtained by peak fitting. The process for fitting has been presented in Chapter \ref{Chap:peakfitting}. Peak fitting is performed from the ``Peak List'' data tree subentry of a histogram and results from the fitting are found as well in the ``Peak List'' data tree entry. When performing peak fitting, it helps to be on the lookout for peaks with shoulders from a nearly overlapped reflection. Inclusion of the second peak, even when the peaks have minimal separation will provide peak positions that are closer to the reflection positions. When fit as a single peak, the reported position will be something of a weighted average of the overlapped reflections. 

The peaks that will be used for autoindexing are found in a separate data tree subentry, ``Index Peak List.'' The reason that this second list is present is that this allows for selection of peaks and possible editing as well as tracking the assignment of \textit{hkl} indices. Use the Operations/``Load/Reload'' command to copy the peak positions from the  ``Peak List'' to the ``Index Peak List''. Note that each peak in this list has a ``Use'' flag that can be turned off to omit the peak position from the indexing process. 

You want to use some care with the choice and number of peak positions. The lowest Q peaks are typically key for indexing as they help define the length of the axes. Autoindexing typically requires a minimum of 20 peak positions but inclusion of more than $\sim 30$ peaks usually makes indexing more difficult. It is usually better to leave out ``suspect'' peaks than include them. What should make you suspicious about a peak? Certainly if you see a peak that is much narrower than the others, this could be arising from a single crystal that is not representative of the sample. Significantly broader peaks could be from a more poorly crystallized impurity or could be from more than one partly overlapped reflection. In this latter case, since the position matches neither reflection, including that will likely degrade the indexing process. 

Once you have populated the ``Index Peak List'' and have selected which peaks to use, you can select the ``Unit Cells List'' data tree tab. The upper part of the resulting window has options for autoindexing and the process used is documented in the ``Fit Peaks/Autoindexing in GSAS-II'' tutorial (\url{https://advancedphotonsource.github.io/GSAS-II-tutorials/FitPeaks/Fit\%20Peaks.htm}). While it is possible to attempt indexing in all the Bravais lattice options at one time, I recommend using a small number at a time. 

Trial unit cells located in the autoindexing process are placed in the lower section of the ``Index Peak List'' window. You will then want to review the located cells visually, to see how well each cell indexes the pattern. When you select a unit cell from this lower list, it is loaded into the middle portion of the window where the reflections generated by that cell can be visualized superimposed on the observed powder diffraction pattern. Note when the  ``Unit Cells List'' data tree entry is selected, the peak positions from the ``Index Peak List'' window are shown as blue lines, when the ``Use'' flag is selected and are not shown otherwise. The lines generated from the selected unit cell are shown as a orange dashed line. Note that the upper data limit is also shown as a red dashed line, but the dashes have double the spacing. Thus, when looking at the plot, lines that have red dashes superimposed on blue lines indicate locations where the observed peaks have been indexed. A blue line that does not have red dashes is an unindexed line -- you will want to make note of these and this can indicate the location of an impurity or that the unit cell is a subcell of one that will index all the lines, but could also be just a cell that does not match the pattern. Locations that have only the orange dashed lines indicate generated reflection positions that do not have associated peaks. Sometimes you will notice small peaks in those locations once you know where to look. An indexed location without intensity is not in itself a serious problem. It could be an extinct reflection or may just be allowed but very weak, but too many unobserved reflections is probably an indication that the cell being evaluated is a supercell of the correct unit cell. 

Note the keep column in the unit cells list. If you find a cell that shows promise, check this. Then if any subsequent searches are made, those entries will be retained while all others are deleted to make room for the new search results. 

\section{Manual indexing}
Manual indexing can be done by entering cell lengths and angles and you can then manually change lengths using the up and down arrows next to the length or angle to increment or decrement the value. This can be a hopeless task with a complex pattern and a low symmetry cell, but can work well for indexing cubic, hexagonal or tetragonal cells, where if two peaks can be indexed as a \textit{h}00 and a 00\textit{l} reflection, the lattice has then been defined. On occasion more complex cells can be indexed if there are low index peaks along each axis. I have also had success with this type of indexing where electron microscopy can identify a 2D unit cell and only a third lattice dimension and possibly an angle needs to be determined. 

Note that autoindexing requires a minimum of 20 peaks, so for minor phases, this is not usually possible. Database searching, for example with the ICDD or COD databases may help identify the phase if already known. For new phases, manual indexing may be the only hope. 

\section{Testing for supercells}

When I have been able to index most of the peaks in a powder diffraction pattern, but there are remain a small number of unindexed peaks, I will immediately think that an impurity is present.  However, one should also consider that these peaks arise from the presence of lowered symmetry that either removes a  reflection extinction condition or expands the unit cell. There are two ways that one can look for this. 

One way is to generate subgroups, using the ``Run Subgroups'' command in the ``Cell Index/Refine'' menu. This will generate a list of unit cells and space groups that can be considered, as was done with autoindexing results. This is described in more detail in \S\ref{ref:subgroups}.

The other approach is simply to look if doubling (or very rarely) tripling one or more unit cell lengths will index the previously unindexed lines. As has been discussed in \S\ref{ref:supercellindexing}, you can quickly expand a unit cell length(s) that are specified on the ``Unit Cells List'' data window by placing a \texttt{2*} before or a \texttt{*2} after and the value will be doubled. This will allow you to see quickly if an expanded lattice will index these extra lines. If so, looking more carefully at subgroups may then be helpful. 

\section{Electron diffraction}

While I would not have written this some years back, now when faced with a tricky problem indexing a powder diffraction pattern, I would today ask if one should even try. A variant of electron microscopy known as MicroED has advanced tremendously, and for many materials, not only can electron diffraction be used with a powder specimens to obtain the unit cell constants, but also to collect a diffraction data set that can be used to solve the structure. Once that is done, there is still great value in further refining the structure against x-ray and/or neutron data, as electron diffraction has significant problems with obtaining accurate intensity values for structure factors due to dynamical electron diffraction, but this is not the case with x-rays and neutrons. Thus, there are significant advantages in using electron diffraction on unknown materials in combination with x-rays. There are instruments designed specifically for measuring electron diffraction, but more some general purpose many transmission electron microscopes (TEM) can also be used. These are very expensive instruments, but some institutions make them available for external use and at other locations collaboration with the microscopists is welcomed. 
